\section{Command reductions and a complete proof compiler}
\label{sec:command-reductions}

Write a command store as $(t,w,d,h)$, where $w$ is a Cantor-coded stack.
For clarity, $a::w$ denotes one stack item, including the zero-tail
convention $0::0=0$. Retain the axiom command unchanged. Its generated
stack has the form
\[
 L_0::L_1::\cdots::L_7::t::w,
\]
where $L_0,\ldots,L_5$ are the six heads in reverse push order,
$L_6$ is the restricted propositional head, and
$L_7=\mathsf U(d,h)$. Every $L_i$ differs from code $1$: the heads
and propositional axiom are implication codes, and the last item has
quantifier tag $2$. The axiom command clears $d$ and $h$.

\subsection{Two independent simplifications}

The original Shift halts when $t+h=1$ and otherwise maps
\[
 (t,a::w,d,h)\longmapsto(h,w,d,a).
\]
The simpler Shift tests only $h=1$, with the same nonhalting update.
Its soundness invariant needs only a proof of the formula in $h$.

The original successful MP command, with $b=d+x$ and
$\mathsf I(t,b)=h$, returns $(b,0::w,b,0)$. A stronger simplification
returns the consequent directly in the other live formula register:
\begin{equation}\label{eq:new-mp-audit}
 (t,w,d,h)\overset{\mathsf{MP}(x)}\longmapsto
 \begin{cases}
 (0,w,0,d+x),&\mathsf I(t,d+x)=h,\\
 (0,w,d+x,0),&\text{otherwise}.
 \end{cases}
\end{equation}
This removes the successful branch's zero push and its copying of the
consequent back to $t$ while retaining $d$. The failed branch is unchanged.
If both live formula registers contain derivable formulas, the successful
branch puts a derivable consequent into $h$ and truth into $t$; the failed
branch puts truth in both. Hence the same syntactic proof invariant applies.
The two simplifications can be used together.

At the register-program level, the original Shift preparation first
transfers $T$ to the scratch register $X$, then transfers $H$ to both
$T$ and $X$, then transfers $X$ to $H$. Its effect is
$(T,H,X)=(h,t+h,0)$. Replace only the first transfer by the existing
clear-$T$ loop. The new effect is $(T,H,X)=(h,h,0)$, so the unchanged
one-test guard now tests the incoming $h$ alone. Counting source
instructions, the three loops take respectively $2t+1$, $4h+1$, and
$3h+1$ steps, a total of $2t+7h+3$. The guard takes eight steps when
$h=1$ and $2h+7$ when $h\ne1$; hence the halting branch takes
$2t+18$ steps and the nonhalting branch reaches stack decoding in
$2t+9h+10$ steps. These counts concern source instructions and require
the corresponding source-address changes; they are not tape-step counts.

For~\eqref{eq:new-mp-audit}, after a successful comparison the live
registers satisfy $T=H=X=0$, $D=b$. The existing destructive-transfer
loop from $D$ to $H$ has $3b+1$ source steps and leaves
$(T,H,D,X)=(0,b,0,0)$. It is therefore sufficient in place of the
zero-push and the preserving transfer from $D$ to $T$.

\subsection{Unchanged auxiliary words}

Use the original words
\[
\begin{split}
 \mathsf{PUSH}&=\mathsf{AX}(0,0,0)\;\mathsf{SHIFT}^{8}\;\mathsf{MP}(0),\\
 \mathsf{DROP}&=\mathsf{MP}(1)\;\mathsf{AX}(1,0,0)\;
                   \mathsf{SHIFT}^{9}\;\mathsf{MP}(0).
\end{split}
\]
Their exact contracts remain
\[
 \mathsf{PUSH}(t,w,d,h)=(0,t::w,0,0),\qquad
 \mathsf{DROP}(t,w,0,0)=(0,w,0,0).
\]
Every MP instruction in these words takes its failed branch. For PUSH,
after eight shifts the value in $h$ is a quantifier code, whereas an
implication code is divisible by four, so equality cannot hold. For DROP,
the first test compares the positive code $\mathsf I(t,1)$ with zero;
the final test compares $\mathsf I(\mathsf U(1,0),0)$ with zero.
The offset word $\mathsf{MP}(u)$ for $u>0$, applied to $(0,w,0,0)$,
also fails and establishes $(0,w,u,0)$. Thus these contracts do not depend
on the old successful MP branch.

The eight- and nine-shift identities following AX are unchanged even with
the new Shift. The former exposes $L_6$ and $L_7$ while retaining the saved
top in the stack; the latter gives
\begin{equation}\label{eq:ax-nine-audit}
 \mathsf{AX}(c,x,y)\;\mathsf{SHIFT}^{9}(t,w,d,h)
       =(\mathsf U(d,h),w,0,t).
\end{equation}
All nine tests examine either zero or one of the non-one generated heads.

The original word $\mathsf{POP}=\mathsf{AX}(0,0,0)\mathsf{SHIFT}^{10}$
now has the contract
\[
 \mathsf{POP}(t,a::w,d,h)=
 \begin{cases}
 \text{halt},&t=1,\\
 (t,w,0,a),&t\ne1.
 \end{cases}
\]
This distinction matters: POP is no longer unconditionally total under
the $h$-only halt test. It does not affect the axiom selectors because
their selected heads and gates differ from $1$. The original
generalization word is unchanged: it uses PUSH, an offset, one Shift,
then AX and nine Shifts, and returns
$(\mathsf U(i,a),w,0,0)$ from $(a,w,0,0)$.

\subsection{Compiling modus ponens with a zero below the premises}

Suppose $p$ proves $A$ of code $a$ and $q$ proves $A\to B$ of code
$\mathsf I(a,b)$. Define the new compiler recursively by retaining
the original axiom selectors and generalization clause, and replacing
its MP clause with
\begin{equation}\label{eq:new-mp-word-audit}
 \mathsf{DROP}\;\mathsf{PUSH}\;C(q)\;\mathsf{PUSH}\;
 C(p)\;\mathsf{PUSH}\;\mathsf{SHIFT}^{2}\;
 \mathsf{MP}(b)\;\mathsf{SHIFT}.
\end{equation}
The initial DROP and PUSH insert a zero below both premises. Starting
with $(t,w,0,0)$, unless a recursive call already halts, the states at
the significant boundaries are
\[
\begin{array}{c|c}
 \text{after the first DROP and PUSH}&(0,0::w,0,0)\\
 \text{after }C(q)\text{ and PUSH}&(0,\mathsf I(a,b)::0::w,0,0)\\
 \text{after }C(p)\text{ and PUSH}&(0,a::\mathsf I(a,b)::0::w,0,0)\\
 \text{after two Shifts, if }a\ne1&(a,0::w,0,\mathsf I(a,b))\\
 \text{after MP}(b)&(0,0::w,0,b)\\
 \text{after the final Shift, if }b\ne1&(b,w,0,0).
\end{array}
\]
The second Shift halts if $a=1$, and the last Shift halts if $b=1$.
Both conclusions hold for either halt test: at these tests $t=0$.
The compiler therefore preserves an arbitrary incoming stack and returns
the derived formula with cleared caches, or halts earlier.

Induction on the finite derivation proves completeness. For a derivation
of code $1$, no axiom or generalization can be its final constructor,
because their conclusion codes differ from $1$. Its last constructor
is MP, whose final Shift halts. Thus the compiled word itself halts.
Prepending one AX resets arbitrary initial caches, exactly as in the
upstream compiler. The construction is primitive recursive and uses
only fixed words and recursive traversal of a finite derivation; its
correctness proof is an arithmetical induction available in PA.

\subsection{A transport proof for the Shift change alone}

There is also a direct way to preserve the original refutation compiler
when only Shift changes. Given an old halting word $W$ and its initial
store, evaluate the finite word to locate its first halting Shift. Let
$U$ be the preceding prefix and $(t,w,d,h)$ its store. Evaluation is
bounded by the word length, so this transformation is primitive recursive.
The new execution of $U$ either already halts, or reaches that identical
store: every nonhalting update is unchanged. If $h=1$, append one Shift.
Otherwise the old condition $t+h=1$ forces $t=1,h=0$; append
\[
 \mathsf{AX}(0,0,0)\;\mathsf{SHIFT}^{10}.
\]
After nine Shifts, equation~\eqref{eq:ax-nine-audit} puts the saved $1$
in $h$, and the tenth Shift halts. This supplies an internal-PA
refutation-to-new-word map without relying on the old final Shift alone.
It must be applied to the prefix before that old halt: simply appending
the postlude after the entire old word can discard the only exposed $1$.

\subsection{A weaker MP option}

If only the zero push is omitted and successful MP instead returns
$(b,w,b,0)$, the same initial DROP and PUSH in the MP compiler clause
provide the missing zero. Keep the old MP suffix after the two premise
pushes: two Shifts, MP$(b)$, AX$(b,0,0)$ and ten Shifts. It then returns
$(b,w,0,0)$ or halts early. The stronger branch
\eqref{eq:new-mp-audit} has the shorter proof-word suffix above and can
also shorten the register program by replacing the preserving copy with
a destructive transfer into $h$.
